\chapter{Locating the speaker}\label{ch:loc}

\section{Why a likelihood grid}

Classic resection solves for the position from exact angles to known points. Spoken reports are not
exact. ``Church at 2 o'clock'' is $\pm$15$^\circ$ at best, the speaker doesn't say \emph{which}
church, sometimes he sees no landmark at all, and some statements are plain wrong. So we don't solve
equations. We score every place on a grid against every observation, and let wrong matches fail on
their own.

\begin{itemize}
  \item \textbf{Grid:} 25 m cells over the 4.4 $\times$ 3.9 km area: 176 $\times$ 155 = 27,280 cells.
  \item \textbf{Heading:} the direction the speaker faces is unknown, so each cell also carries
        72 headings (5$^\circ$ steps). A clock direction only means something relative to the facing
        direction.
  \item \textbf{Output:} a probability map over cells, after summing out the heading.
\end{itemize}

\nanofig{resection}{Resection from spoken clock directions. With the facing direction unknown, two landmarks
leave a family of candidate positions; a third landmark, the sun or a compass bearing picks one.}
{Clean geometric diagram, top-down view, plain white background, flat vector. A small town map in light grey with a
castle keep icon at top-left and a church icon at top-right and a water tower icon at bottom-right. A soldier
symbol (small circle with an arrow showing facing direction) at the centre. From the soldier, a cyan line to the
keep labelled ``11 o'clock'', an amber line to the church labelled ``2 o'clock''. A faint clock face drawn around
the soldier with 12 aligned to his facing arrow. A dashed cyan circle arc passing through the keep, the church and
the soldier, labelled ``positions that fit two angles''. A small sun icon at the upper right with a thin arrow labelled
``sun at 1 o'clock''. Title at top: ``Pothenot's problem, spoken''. Exactly these labels. Colour palette (use only these): navy \#0D1110, cyan \#1AA7C4, amber \#E09A2B, red \#E5484D, green \#3FA877, grey \#8B9A94, plus white for text on dark shapes.}

\section{The model}

For a cell $p$ and a heading $\theta$, each constraint $i$ gives a likelihood $\ell_i(p,\theta)$. They
are combined with an outlier floor $\varepsilon$, so a single wrong statement can't eliminate the true
position:
\begin{equation}
  L(p,\theta) \;=\; \prod_i \Big[\varepsilon + (1-\varepsilon)\,\ell_i(p,\theta)\Big],
  \qquad \varepsilon = 0.01 .
\end{equation}
The position posterior sums over headings:
\begin{equation}
  P(p) \;\propto\; \sum_{\theta} L(p,\theta).
\end{equation}

\paragraph{Landmark bearings (clock directions).} The speaker said landmark type $T_i$ (optionally a
name) at relative bearing $\beta_i = 30^\circ \times$ clock hour. Let $C_i$ be the candidate landmarks
of that type and $b(p,c)$ the true bearing from $p$ to $c$. Then
\begin{equation}
  \ell_i(p,\theta) \;=\; \max_{c \in C_i}\; w_c(p)\,
  \exp\!\left(-\frac{\operatorname{wrap}\big(b(p,c) - \theta - \beta_i\big)^2}{2\sigma_i^2}\right).
\end{equation}
The \emph{max} does the matching: the wrong church simply doesn't fit. The weight $w_c(p)$ says whether
$c$ could be seen from $p$: distance, buildings and terrain from the LiDAR model, and trees
(Section~\ref{sec:3d}). If a distance was said, a log-normal term with $\sigma = 50\,\%$ multiplies in. In code
the max over candidates is computed per heading bin, then spread by the angular noise with a
max-convolution. Each bearing costs one vectorised pass over cells $\times$ candidates.

\paragraph{Compass bearings.} A compass gives the absolute bearing $\alpha_i$ (magnetic, corrected by
+2$^\circ$ declination for Paris). The facing direction disappears:
\begin{equation}
  \ell_i(p) = \max_{c \in C_i} w_c(p)\,\exp\!\left(-\frac{\operatorname{wrap}(b(p,c) - \alpha_i)^2}{2\sigma^2}\right),
  \qquad \sigma = 6^\circ .
\end{equation}
This is the clock-bearing likelihood evaluated at $\beta=0$ and read at heading column $\theta=\alpha_i$,
so it reuses the same cached computation.

\paragraph{The sun as a compass.} \code{sun.py} computes the sun's azimuth $a_\odot$ and elevation from the
time and the centre of the area (NOAA formulas). ``Sun at my $k$ o'clock'' constrains only the heading:
$\ell(\theta) = \exp(-\operatorname{wrap}(a_\odot - \theta - 30k)^2/2\sigma^2)$ with $\sigma = 15^\circ$.
It's ignored when the sun is below the horizon or the speaker says it's overcast. With the heading fixed,
two landmarks are enough.

\paragraph{Phone compass and gyroscope.} ``My 12 o'clock'' taken with the phone flat gives a heading
constraint at $\pm$25$^\circ$ (phone compasses held still). Aim \& tap gives the elevation angle of a
landmark's top and the turn between two tops (Section~\ref{sec:3d}).

\paragraph{Other clues.} ``Can't see a church'' lowers cells that would see one. ``Under trees, on a dirt
track'' uses the forest grid. Sounds, the kind of space, wall echoes and the noise level come from the
acoustic chapter, at weight 0.7. They narrow the \emph{kind} of place and never override landmarks. A last
GPS fix before jamming bounds the area to walking range: 1.4 m/s since the fix.

\paragraph{Walking prior.} When a guided walker reports new landmarks, the old observations are dropped
(they were seen from elsewhere) and replaced by a Gaussian prior around the dead-reckoning position, with
a much lower outlier floor ($10^{-4}$). The walker followed our route, so being far from it is very unlikely.

\section{From a map to an answer}

\begin{itemize}
  \item \textbf{Areas:} connected regions of high probability, up to three blobs, each with its mass, a
        1-$\sigma$/2-$\sigma$ ellipse from the weighted covariance, and $r_{80}$, the radius holding 80 \% of it.
  \item \textbf{Distinct spots A/B/C:} peaks at least 300 m apart with their share. They appear on the map
        and in the read-back (``Also possible: area B, 900 metres west''), and question goals use them.
  \item \textbf{State:} \emph{ambiguous} when the main area holds less than 45 \% (or two big areas compete),
        \emph{confirmed} when the radius is under 75 m, \emph{rough} otherwise.
  \item \textbf{Read-back radius:} rounded to speakable numbers (``within about 300 metres'').
\end{itemize}

\nanofig{grid}{The likelihood grid: each observation is a layer over cells and headings; their product, summed over
headings, gives the position map with its areas A and B.}
{Exploded isometric technical illustration, 16:9, plain white background, flat vector. Four semi-transparent stacked
square layers floating above a base map of a small town. Layer 1 labelled ``keep at 11 o'clock'' shows a curved
cyan band. Layer 2 labelled ``church, bearing 240°'' shows a straight amber wedge. Layer 3 labelled ``sun at
1 o'clock'' shows a uniform yellow tint with a small heading dial. Layer 4 labelled ``can't see the water tower''
shows a grey hole. At the bottom, the base map shows the product: a bright cyan blob labelled ``A 72\%'' and a
smaller blob labelled ``B 21\%'' with a dashed ellipse around A. A side axis labelled ``72 headings'' on the left.
Exactly these labels. Colour palette (use only these): navy \#0D1110, cyan \#1AA7C4, amber \#E09A2B, red \#E5484D, green \#3FA877, grey \#8B9A94, plus white for text on dark shapes.}

\section{Seeing in 3D: LiDAR line of sight and Aim \& tap}\label{sec:3d}

\paragraph{Who can see what.} The early model assumed a landmark is visible when it's close enough. That's
false in a city: from most street spots the keep is hidden behind buildings. \code{viewshed.py} builds a 2 m
terrain and surface model from IGN LiDAR HD. For each of the 304 landmarks with a top (keep, towers,
churches, schools, stations) and each cell in range, it traces the sight line from eye height to the top.
Buildings and terrain block it; trees let a fraction through, $\exp(-d_{\text{wood}}/100\,\text{m})$. The cache
takes about 35 s to build once. From the streets, the keep's top turns out to be visible from only about
11 \% of spots within 1.5 km, mostly within 200 m. The old rule assumed about 40 \%.

\paragraph{Aim \& tap.} The soldier puts the top edge of the phone on the top of each landmark he sees,
taps, and names them in order.
\begin{itemize}
  \item The \textbf{tilt} is the elevation angle $e$ of the top. With the LiDAR top altitude $z_t$ and the
        cell's ground altitude $z_g$, the horizontal distance follows from trigonometric levelling,
        $d = (z_t - z_g - 1.5\,\text{m}) / \tan e$, at $\pm$2$^\circ$.
  \item The \textbf{turn} between two taps (gyroscope, no compass) is the angle between two landmarks seen from
        the spot. That is Pothenot's problem in its exact form, with no heading needed, at $\pm$2.5$^\circ$.
\end{itemize}

\nanofig{aimtap}{Aim \& tap: the phone's tilt to a landmark's top gives its distance; the turn between two taps gives the
angle between landmarks.}
{Technical side-and-top dual view illustration, 16:9, flat vector, plain white background. Left half, side view: a
soldier holding a phone at eye level, the phone's top edge pointing at the top of a tall castle keep 300 m away;
a dotted sight line, an angle arc labelled ``tilt e'', the keep height labelled ``53 m (LiDAR)'', ground line,
distance arrow labelled ``d = (z\_top − z\_ground − 1.5 m) / tan e''. Right half, top view: the same soldier dot with
two sight lines to a keep icon and a church icon, an angle arc between them labelled ``turn between taps''. A
small gyroscope icon. Exactly these labels. Colour palette (use only these): navy \#0D1110, cyan \#1AA7C4, amber \#E09A2B, red \#E5484D, green \#3FA877, grey \#8B9A94, plus white for text on dark shapes.}

\section{In the woods}

Under trees, distant landmarks disappear, the sun often does too, and there are no walls to echo.
\code{forest.py} gives each cell its share of woods, open ground and water, the distance to the nearest
paved and unpaved path, depth in the wood, distance to water and fences. It uses 36 line-of-sight profiles
per cell. Ground clues (``under trees'', ``dirt track'', ``I can see a pond'') are said or asked, and so is the
direction of road traffic heard through the trees.

\section{Results}

\begin{keyresult}[Ask-back loop, spots where at least two landmarks are truly visible]
Simulated speakers see the world through the 3D model and answer truthfully. Up to 3 questions: median error
\textbf{43 m}, 80th percentile 190 m. The spoken radius goes from about 720 m to 80 m. The truth falls inside it
in 73 \% of trials, so the radius is a little too confident. (\code{scripts/eval\_loop.py})
\end{keyresult}

\begin{figure}[htbp]\centering
\begin{tikzpicture}
\begin{axis}[ybar, width=0.95\linewidth, height=6cm, bar width=13pt, ymin=0,
  ylabel={Median error (m)}, symbolic x coords={2 landmarks, 3 landmarks},
  xtick=data, legend style={at={(0.5,1.03)},anchor=south,legend columns=-1},
  nodes near coords, enlarge x limits=0.35, every node near coord/.append style={font=\footnotesize}]
\addplot[fill=greyN!60] coordinates {(2 landmarks,250) (3 landmarks,97)};
\addplot[fill=cyanN] coordinates {(2 landmarks,71) (3 landmarks,59)};
\legend{Clock directions (heading unknown), Compass bearings}
\end{axis}
\end{tikzpicture}
\caption{Compass bearings against clock directions, 30 random spots, named landmarks 80--900 m away and more than
35$^\circ$ apart. Within 75 m with 2 landmarks: 57 \% against 13 \%.}\label{fig:compasschart}
\end{figure}

\begin{table}[htbp]\centering\small
\caption{3D line of sight and Aim \& tap, 60 simulated casualties per area at spots where at least one landmark
top is truly in sight. Same spots and random draws in every row (\code{scripts/eval\_geomatics.py --seen}).}\label{tab:geomatics}
\begin{tabular}{lrrrr}
\toprule
& \multicolumn{2}{c}{Streets} & \multicolumn{2}{c}{Woods}\\
\cmidrule(lr){2-3}\cmidrule(lr){4-5}
& median error & within 50 m & median error & within 50 m\\
\midrule
Before (distance + canopy rule) & 417 m & 7 \% & 492 m & 5 \%\\
+ 3D line of sight & 190 m & 23 \% & 427 m & 8 \%\\
+ Aim \& tap (tilt, turn) & \textbf{75 m} & \textbf{45 \%} & 469 m & 10 \%\\
\bottomrule
\end{tabular}
\end{table}

\begin{table}[htbp]\centering\small
\caption{In the woods, 60 simulated casualties in wooded cells, one report then up to 3 questions
(\code{scripts/eval\_forest.py}).}\label{tab:forest}
\begin{tabular}{lrrr}
\toprule
& median error & within 75 m & within 150 m\\
\midrule
Before & 1,123 m & 3 \% & 5 \%\\
+ canopy-aware visibility & 1,155 m & 0 \% & 3 \%\\
+ ``in the woods, dirt track'', ground and traffic questions & 473 m & 10 \% & 23 \%\\
+ phone compass & 458 m & 12 \% & 22 \%\\
\bottomrule
\end{tabular}
\end{table}

\begin{caveat}
\begin{itemize}
  \item In the streets, 57 \% of random spots see \emph{no} landmark top at all, and 73 \% in the woods. There only
        questions and ground or traffic clues remain, and the median error is around 850 m (streets) and 1.5 km
        (woods). Our earlier city numbers assumed landmarks visible whenever close, and were too optimistic.
  \item Canopy-aware visibility doesn't improve the forest numbers by itself; it removes wrong conclusions such as
        ``can't see the keep, so you're far from it''.
  \item The 3D model plays reality in these tests and is also what the localiser uses. Leaves in summer, trucks,
        scaffolding and landmark points placed off the structure will make real results worse.
  \item The map covers the Vincennes sector only. Another area needs its OSM, IGN and LiDAR data fetched.
\end{itemize}
\end{caveat}
